\section{Worked cases, II}

\subsection{Case 3: All Souls falling on a Sunday}

\begin{casebox}{Case 3. Sunday 2 November and Monday 3 November 2036}

\textbf{Calendar facts.} Easter 2036 falls on 13 April, Pentecost on 1 June,
Trinity Sunday --- the First Sunday after Pentecost --- on 8 June, and Advent
Sunday on 30 November. Counting from Trinity, 2 November 2036 is the
twenty-second Sunday after Pentecost. Christ the King, ``dominica ultima
mensis octobris'' (\rg{17}\,d), fell the previous week on 26 October. All
Saints, 1 November, is the Saturday.

\begin{steps}
\item \textbf{Day, 2 November.} Two first-order competitors: the Twenty-second
Sunday after Pentecost, a second-class Sunday (\rg{12}), and the Commemoration
of All the Faithful Departed, a first-class liturgical day.

\item \textbf{Precedence.} This is the one place where the table hands the
lower row the victory. Row 8 reads ``Commemoratio omnium Fidelium
defunctorum, quae tamen locum cedit dominicae occurrenti,'' and \rg{16}\,b
states it from the other side: ``dominica II classis praefertur Commemorationi
omnium Fidelium defunctorum.'' \textbf{The Sunday takes the day.} Note that
row 8 does not merely yield to a first-class Sunday; it yields to \emph{any}
occurring Sunday, which is why an ordinary green Sunday after Pentecost
defeats a first-class day.

\item \textbf{Impediment.} All Souls is not a feast, so \rg{95}'s right of
transfer, which belongs to first-class feasts, does not reach it. \rg{96}\,b
supplies the rule directly: ``Commemoratio omnium Fidelium defunctorum, quando
occurrit cum dominica, transfertur, tamquam in sedem propriam, in feriam II
sequentem.'' \textbf{Monday 3 November 2036.} The Missal prints the same
result in its own heading at the place: \latin{Die 2 novembris vel, si in
dominicam inciderit, die 3 sequenti}. Nothing of All Souls remains on the
Sunday: there is no commemoration of it, because \rg{111}\,b admits on a
second-class Sunday only the commemoration of a second-class feast.

\item \textbf{Mass category, Sunday 2 November.} The Mass of the Office.
Second-class votives are barred on Sundays for the \latin{Pro Sponsis} Mass
(\rgm{379}) and in general are permitted on second-class days (\rgm{341}), so
a second-class votive is not excluded by the day's class as such; but a
second-class Requiem is, since \rgm{411}\,b bars them whenever ``dies
liturgicus I classis aut dominica quaevis'' occurs, and a third-class Requiem
is barred by \rgm{416}. The funeral Mass is a separate and contested question:
see Section~10.

\item \textbf{Assembly, Sunday 2 November.}
\afield{Formulary}{The Twenty-second Sunday after Pentecost.}
\afield{Colour}{Green (\rg{127}\,b).}
\afield{Gloria}{Said --- Te Deum on second-class Sundays outside the
Septuagesima group (\rgb{237}\,b).}
\afield{Orations}{One. Nothing is admitted alongside it.}
\afield{Creed}{Said (\rgm{475}\,a).}
\afield{Preface}{Of the Most Holy Trinity, as the Preface \latin{de Tempore}
for every second-class Sunday outside Christmastide and Paschaltide
(\rgm{494}\,b).}
\afield{End}{\latin{Ite, missa est}; blessing; Last Gospel of St John.}

\item \textbf{Monday 3 November: the day.} All Souls, first class, in its
proper seat. The Missal's rubric at the head of the day:

\begin{quote}\itshape\small
Hac die quivis sacerdos tres Missas celebrare potest. \dots{} Qui unam
dumtaxat Missam celebrat, primam legit: eandem adhibet qui Missam cum cantu
celebrat, facta ei potestate anticipandae secundae ac tertiae. Cum quis tres
Missas sine intermissione celebrat, sequentiam dicere debet tantum in Missa
principali, secus in prima Missa; in ceteris Missis, nisi sint in cantu, eam
omittere potest.
\end{quote}

\rgm{403}--\rgm{404} say the same and add that a priest saying several sung
Masses in different churches must always use the first, and that where several
sung Masses are said in one church the first, second and third are used in
order. \rgm{295} makes the first Mass the conventual Mass, ``et chorales illi
soli interesse tenentur.''

\item \textbf{Monday 3 November: assembly.}
\afield{Class}{A Mass of the dead of the first class (\rgm{402}\,a).}
\afield{Colour}{Black (\rg{132}\,b), with the exception of \rg{130}\,d: violet
where the Masses of All Souls are celebrated during exposition for the Forty
Hours.}
\afield{Prayers at the foot}{Psalm \latin{Iudica me} omitted (\rgm{425}\,b);
Confiteor as usual.}
\afield{Introit}{\latin{Requiem aeternam}, without \latin{Gloria Patri}
(\rgm{428}); the Missal directs the antiphon repeated ``absolute \dots{} usque
ad psalmum.''}
\afield{Gloria}{Omitted (\rgm{432}\,d; and Te Deum is omitted in the Office of
the dead, \rgb{238}\,d).}
\afield{Orations}{One, and it must be for the dead: ``in Missis defunctorum
prohibetur quaevis oratio, quae non est pro defunctis'' (\rgm{398}\,c), and
``in Missis defunctorum nulla fit commemoratio de Officio diei currentis''
(\rgm{391}). \rgm{394} directs that where the Mass is for a deceased pope,
cardinal, bishop or priest, the first Mass of the day is used with the proper
orations found among the \latin{orationes diversae pro defunctis}.}
\afield{Sequence}{\latin{Dies irae} is obligatory: \rgm{399}\,a confines the
obligation to first-class Masses of the dead, and this is one --- with the
Missal's own relief for a priest saying three Masses without interruption.}
\afield{Creed}{Omitted (\rgm{476}\,f).}
\afield{Preface}{Of the dead (\rgm{499}).}
\afield{Pax}{Not given; the Missal's \work{Ordo Missae} directs at the place
``In Missis defunctorum non datur pax, neque dicitur praecedens oratio.''}
\afield{End}{\latin{Requiescant in pace}, answered \latin{Amen}
(\rgm{507}\,c); \emph{no blessing} (\rgm{508}); and the Last Gospel
\emph{omitted if the absolution at the tomb follows} (\rgm{510}\,e), said if it
does not. \rgm{401}\,a requires the absolution after the funeral Mass and
merely permits it after the others, so the answer differs Mass by Mass on the
same morning.}
\afield{Also barred that day}{The votive Mass \latin{Pro Sponsis}, its
commemoration in the Mass of the day, and the nuptial blessing within Mass ---
all three (\rgm{381}\,d). And \rgm{352} governs the Blessed Sacrament if it is
exposed for the Forty Hours.}

\item \textbf{The consequence nobody expects.} \rgm{422} allows all Masses
applied for the dead to be said as Requiems ``intra octiduum a die
Commemorationis omnium Fidelium defunctorum inclusive computandum.'' In 2036
the day of the Commemoration is 3 November, so on the face of the words the
octave runs from 3 to 10 November rather than from 2 to 9. This is a necessary
inference from the rubric's own wording, not a rule stated in terms; we mark it
as an inference and would defer to a competent Ordo.
\end{steps}
\end{casebox}

\subsection{Case 4: a First Friday votive on a day that is not free}

\begin{casebox}{Case 4. Friday 3 December 2027 --- the same day as Case 1b}

\begin{steps}
\item \textbf{Day.} Established in Case 1b: St Francis Xavier, third-class
feast of the universal calendar, takes the day from the Advent feria, which is
commemorated.

\item \textbf{The request.} The parish keeps the First Friday devotions and
asks for the Mass of the Sacred Heart.

\item \textbf{Mass category --- the test that actually decides.} \rgm{385}\,b
permits \emph{two} Masses of the Sacred Heart ``prima feria sexta cuiusque
mensis, in ecclesiis et oratoriis in quibus peculiaria pietatis exercitia in
honorem eiusdem Ss.mi Cordis, eo die, peraguntur.'' Three conditions, all
independent of the calendar: it must be the first Friday; the exercises must be
performed; and they must be performed \emph{in this church, that day}. The
permission is not the priest's, and it does not travel with him.

The class test then applies. This is a votive Mass of the third class, and
\rgm{384} admits third-class votives ``diebus liturgicis III et IV classis.''
Today is third class. \textbf{Admitted.} Had St Francis Xavier been a
second-class Patron in the local calendar, or had 3 December fallen on a
second-class Sunday, the same request would have been refused --- and
\rgm{386}\,d adds that when third-class votives are prohibited, ``non
commemorantur in Missa diei.'' There is no consolation collect.

Two further bars must be checked and are clear here: \rgm{317} forbids a
votive of a mystery of the Lord when a first- or second-class day occurs whose
Office is of the same Person, which is not the case; and \rgm{326}\,a forbids
every votive in a church having only one Mass when a conventual obligation
presses that no other priest can discharge.

\item \textbf{Assembly.}
\afield{Formulary}{\rgm{314}: ``Pro Missa votiva de mysteriis Domini sumitur
Missa de respectivo festo'' --- that is, the Mass \latin{Cogitationes Cordis
eius} of the feast of the Sacred Heart, printed at the Friday after the Second
Sunday after Pentecost.}
\afield{Colour}{White. \rg{121}\,a requires white for votive Masses answering
the feasts of \rg{120}, which include feasts of the Lord ``exceptis festis de
mysteriis et instrumentis Passionis.'' The Sacred Heart is not among the
excepted.}
\afield{Gloria}{\textbf{Said} --- on a ferial-coloured day in Advent.
\rgm{431}\,d gives the Gloria to votive Masses of the first, second and third
class ``nisi adhibeatur color violaceus paramentorum,'' and the colour here is
white, not violet.}
\afield{Creed}{Omitted. \rgm{386}\,a: third-class votives are said ``cum
Gloria; sed semper sine Credo.'' \latin{Semper}.}
\afield{Orations}{Three in all, which is the maximum \rgm{435} allows.
\rgm{386}\,b admits two commemorations besides the Mass's own oration. First
the collect of the Sacred Heart; then the commemoration of the Advent feria,
which is privileged (\rg{109}\,e), must be made (\rg{25}), and goes first among
the commemorations because it is of the season (\rg{113}); then the
commemoration of St Francis Xavier, whose Office is the Office of the day.
Secrets and postcommunions in the same number and order (\rgm{480},
\rgm{505}).}
\afield{Sequence}{None --- \rgm{470} omits the sequence in every votive Mass,
and this Mass has none in any case.}
\afield{Preface}{Of the Sacred Heart, ``in Missis festivis et votivis de
Ss.mo Corde Iesu'' (\rgm{491}). Not the Common Preface that the day would
otherwise take (Case 1b), because step one of \rgm{482} is satisfied.}
\afield{If sung}{The solemn tone for the orations, Preface and Pater noster
(\rgm{386}\,c, \rgm{515}\,f). But note \rgm{434}\,a: at a \emph{sung
non-conventual} Mass only one privileged commemoration is admitted, so the
commemoration of St Francis Xavier drops out and only the Advent feria
survives. The Low Mass and the sung Mass of the same votive on the same
morning have different numbers of collects.}
\afield{End}{\latin{Ite, missa est}; blessing; Last Gospel of St John.}
\end{steps}

\textbf{Why this case is instructive.} Nothing in the day's class number told
us whether the votive was admitted; the class number only told us it was
\emph{not excluded}. The operative facts were that the exercises are actually
held in this church today and that the church is not under a pressing
conventual burden. Then, having been admitted, the votive changed the Gloria,
the Preface, the colour and the number of orations --- and the manner of
celebration, sung or read, changed the number again.
\end{casebox}

\subsection{Case 5: assembling a Sunday out of two Sundays}

\begin{casebox}{Case 5. The autumn of 2026 --- the resumed Sundays after
Epiphany}

\textbf{Calendar facts.} Easter 2026 falls on 5 April; Septuagesima on 1
February; Pentecost on 24 May; Trinity Sunday on 31 May; Advent Sunday on 29
November. Epiphany, 6 January, is a Tuesday. The Sundays after Epiphany
actually celebrated in January 2026 are the First (11 January), the Second (18
January) and the Third (25 January); Septuagesima then supervenes, and the
Fourth, Fifth and Sixth after Epiphany are impeded. Counting Trinity as the
First Sunday after Pentecost and 22 November as the last Sunday before Advent,
there are twenty-six Sundays after Pentecost.

\begin{steps}
\item \textbf{The rule.} \rg{18}: ``Dominicae post Epiphaniam quae,
superveniente Septuagesima, impediuntur, transferuntur post dominicam XXIII
post Pentecosten, hoc ordine \dots{} b) si dominicae fuerint 26, dominica XXIV
erit quae inscribitur V post Epiphaniam; et XXV, quae inscribitur VI.'' And the
closing paragraph: ``Ultimo tamen loco semper ponitur ea quae in ordine est
XXIV post Pentecosten, omissis, si opus sit, ceteris, quae aliquando locum
habere non possunt.''

\item \textbf{The result for 2026.} The Twenty-third Sunday after Pentecost is
1 November; the Twenty-fourth slot, 8 November, is filled by the formulary
headed \latin{Dominica quinta quae superfuit post Epiphaniam}; the
Twenty-fifth, 15 November, by \latin{Dominica sexta quae superfuit post
Epiphaniam}; and the Twenty-sixth and last, 22 November, by the formulary
headed \latin{Dominica XXIV et ultima post Pentecosten}, which by the rubric's
final paragraph is always last.

\item \textbf{The Sunday that vanishes.} The Fourth Sunday after Epiphany was
impeded in January and is \emph{not} resumed, because the count of Sundays
after Pentecost, not the count of Sundays actually lost, decides how many are
resumed. \rg{14} has already said that an impeded Sunday is neither anticipated
nor resumed; \rg{18} is the narrow exception, and it does not stretch. The
Fourth Sunday after Epiphany simply does not occur in 2026.

\item \textbf{Assembly --- a Mass made of two Sundays.} \rgm{298}:

\begin{rulebox}{\rgm{298}}
Omnes dominicae, sive I sive II classis, propriam Missam habent. Attamen
dominicae post Epiphaniam, quae transferuntur inter dominicam XXIII et XXIV
post Pentecosten, sumunt antiphonas ad Introitum, ad Offertorium et ad
Communionem, necnon graduale et Alleluia cum suo versu a dominica XXIII post
Pentecosten, retentis orationibus, Epistola et Evangelio propriis.
\end{rulebox}

So on 8 November 2026 the Introit is \latin{Dicit Dominus: Ego cogito
cogitationes pacis} --- the Introit of the Twenty-third Sunday after Pentecost
--- while the collect, Epistle and Gospel are those of the Fifth Sunday after
Epiphany. The Missal saves the celebrant the assembly by printing the composite
formulary in full at its place among the Sundays after Pentecost, Gradual,
Alleluia, Offertory and Communion all taken over.

\afield{Colour}{Green (\rg{127}\,b).}
\afield{Gloria and Creed}{Both said (\rgb{237}\,b; \rgm{475}\,a).}
\afield{Preface}{Of the Most Holy Trinity, as for every second-class Sunday
outside Christmastide and Paschaltide (\rgm{494}\,b); the Missal prints the
direction in the formulary.}
\afield{Orations}{One, unless a second-class feast occurs, in which case one
commemoration of it is admitted and is dropped if a privileged commemoration is
due (\rg{111}\,b).}
\end{steps}

\begin{witnesstext}{English witness, the collect of the Fifth Sunday after
Epiphany (Cummiskey, Philadelphia, 1861).}
Preserve, we beseech thee, O Lord, thy family by thy constant mercy, that we,
who confide solely in the support of thy heavenly grace, may be always defended
by thy protection. Thro'.
\end{witnesstext}

\noindent
That is quoted as an identified nineteenth-century lay hand missal's English,
to show what the retained collect says; the 1861 book is a pre-1955 witness and
its rubrics, second and third collects, and Holy Week are not those of 1962.
The 1962 formulary keeps only the Sunday's own collect.

\textbf{The harder variant, and where it runs out.} In 2035, Easter falls on 25
March --- the same year as Case 2 --- and there are twenty-eight Sundays after
Pentecost, so \rg{18}\,d resumes four Epiphany Sundays at slots XXIV to XXVII.
Slot XXIV is 28 October 2035, which is the last Sunday of October and therefore
the feast of Christ the King, a first-class feast which by \rg{17}\,d takes the
place of the occurring Sunday with no commemoration of it. One of the four
resumed Sundays consequently has nowhere to go.

\begin{substeps}
\item \emph{Reading one.} \rg{18} assigns each resumed Sunday to a numbered
slot. Slot XXIV is occupied by Christ the King, so the Sunday assigned to it
--- the Third after Epiphany --- is impeded, and \rg{14} forbids anticipating or
resuming an impeded Sunday. The Third after Epiphany is lost; the Fourth,
Fifth and Sixth are kept at slots XXV, XXVI and XXVII.
\item \emph{Reading two.} \rg{18}'s ``hoc ordine'' states a sequence rather
than a slot assignment, and its closing clause --- ``omissis, si opus sit,
ceteris, quae aliquando locum habere non possunt'' --- is designed for exactly
this shortage. The Sundays are taken in order over the Sundays actually
available after the Twenty-third, so the Third, Fourth and Fifth after
Epiphany are said on 4, 11 and 18 November and the Sixth is dropped.
\end{substeps}

\noindent
The two readings put a different Mass on three consecutive Sundays. We do not
think the general rubrics settle it, and we do not assert an answer: a
competent Ordo for 2035 does. What the rubrics do settle is that one of the
four is lost and that the Twenty-fourth and last Sunday after Pentecost is said
on 25 November whichever reading is taken.
\end{casebox}

\subsection{Case 6: a funeral on All Souls}

\begin{casebox}{Case 6. A funeral requested for Monday 3 November 2036}

\begin{steps}
\item \textbf{Is it barred?} \rgm{406} lists five prohibitions of the funeral
Mass: on the days under rows 1, 2, 3, 4, 5 and 6 of the precedence table; on
days of precept among the feasts comprised under row 11; on the anniversary of
the dedication and the feast of the Title of the church where the funeral is
held; on the feast of the principal Patron of the town or city; and on the
feast of the Title and of the holy Founder of the Order to which that church
belongs. \textbf{All Souls is row 8.} It is not in the list, and \rgm{393},
which bars every Mass of the dead during exposition, expressly excepts the
Masses of All Souls. The funeral Mass is admitted.

\item \textbf{Which formulary.} \rgm{409} anticipates the case: ``In
Commemoratione omnium Fidelium defunctorum pro Missa exsequiali sumitur prima
Missa diei cum orationibus pro respectivo defuncto in Missa exsequiali
dicendis. Si vero prima Missa celebratur pro Officio diei, pro Missa exsequiali
sumitur secunda, aut denique tertia.'' So the funeral takes the first of the
day's three Masses with the funeral orations substituted --- unless the first
is already being said for the Office of the day, in which case the second, and
failing that the third.

\item \textbf{Which orations.} \rgm{395} assigns the formulary \latin{In die
obitus seu depositionis defuncti} to the funeral Mass of a person who is not a
priest; for a deceased pope, cardinal, bishop or priest, \rgm{394}\,a routes
the celebrant instead to the first Mass of All Souls with the proper orations
from the \latin{orationes diversae pro defunctis}.

\item \textbf{Assembly.} As Case 3, with three differences. The Mass is of the
first class (\rgm{402}\,b), so \latin{Dies irae} is obligatory (\rgm{399}\,a).
The absolution at the tomb ``fieri debet post Missam exsequialem''
(\rgm{401}\,a), so the Last Gospel is omitted (\rgm{510}\,e). And the orations
are those for this deceased person, one only, since \rgm{398}\,a gives every
Mass of the dead a single oration unless an imposed prayer for the dead under
\rgm{458} or a votive prayer for the dead under \rgm{464} may be added --- and
neither may be here, since \rgm{458} confines the imposed prayer for the dead
to fourth-class ferias and to fourth-class read votive and Requiem Masses.

\item \textbf{If the funeral cannot be held that day.} \rgm{408}: whenever the
funeral Mass is prohibited, or for a reasonable cause cannot be celebrated at
the funeral itself, ``transferri potest in proximiorem diem similiter non
impeditum.''
\end{steps}

\textbf{What this case shows.} The intuition that a Requiem must be
inadmissible on a great day of the dead is backwards; the day is a Requiem. The
intuition that a first-class day excludes a funeral is also backwards; the
prohibition list is a list of named rows and named feasts, not a class
threshold --- which is the point argued in the synthesis at the end of
Section~5, and which the contested reading in Section~10 tests from the other
direction.
\end{casebox}
